% Single-file source generated from the modular article.
\documentclass[11pt]{article}
% BEGIN INCLUDED FILE: preamble.tex
\usepackage[T1]{fontenc}
\usepackage[utf8]{inputenc}
\usepackage{lmodern}
\usepackage[letterpaper,margin=0.94in,headheight=14pt]{geometry}
\usepackage{amsmath,amssymb,amsthm,mathtools,mathrsfs}
\usepackage{microtype,booktabs,array,longtable}
\usepackage[dvipsnames]{xcolor}
\usepackage{enumitem,fancyhdr}
\usepackage{hyperref}
\hypersetup{colorlinks=true,linkcolor=MidnightBlue,citecolor=MidnightBlue,urlcolor=MidnightBlue,
 pdftitle={A Cubic-Background Formula for Gamma Normal Forms},pdfauthor={ProveIt research continuation}}
\urlstyle{same}
\pagestyle{fancy}\fancyhf{}
\fancyhead[L]{\small\sffamily Gamma normal forms and spectral identities}
\fancyhead[R]{\small\sffamily ProveIt continuation}
\fancyfoot[C]{\thepage}
\renewcommand{\headrulewidth}{0.3pt}
\setlength{\parindent}{1.2em}\setlength{\parskip}{2pt}
\emergencystretch=1.6em
\allowdisplaybreaks[2]
\numberwithin{equation}{section}
\newtheorem{theorem}{Theorem}[section]
\newtheorem{proposition}[theorem]{Proposition}
\newtheorem{lemma}[theorem]{Lemma}
\newtheorem{corollary}[theorem]{Corollary}
\newtheorem{conjecture}[theorem]{Conjecture}
\theoremstyle{definition}\newtheorem{definition}[theorem]{Definition}
\theoremstyle{remark}\newtheorem{remark}[theorem]{Remark}
\newcommand{\Q}{\mathbb Q}\newcommand{\C}{\mathbb C}
\newcommand{\Rcal}{\mathscr R}
\newcommand{\NF}{\mathsf N}\newcommand{\Prim}{\mathsf H}
\newcommand{\D}{\mathsf D}\newcommand{\Euler}{\mathscr E}
\newcommand{\Acal}{\mathcal A}
\newcommand{\avg}[1]{\left\langle #1\right\rangle}
\newcommand{\dd}{\,\mathrm d}
\newcommand{\wt}{\operatorname{wt}}
\newcommand{\Shi}{\operatorname{Shi}}
\newcommand{\Li}{\operatorname{Li}}
\newcommand{\FP}{\operatorname{FP}}
\newcommand{\fall}[2]{#1^{\underline{#2}}}
\newcommand{\Erem}{\mathcal E_2}
\newcommand{\eval}[1]{\left[#1\right]_{-1}^{1}}

% END INCLUDED FILE: preamble.tex

\begin{document}
\hypersetup{pageanchor=false}
\begin{titlepage}
\setlength{\parindent}{0pt}
\thispagestyle{empty}
{\sffamily\color{MidnightBlue}\large PROVEIT RESEARCH CONTINUATION\par}
\vspace{1cm}
{\huge\bfseries A Cubic-Background Formula\\[5pt]
for Gamma Normal Forms\par}
\vspace{0.5cm}
{\Large Complete mixed polygamma coefficients and\\[4pt]
Stieltjes--polylogarithm endpoint identities\par}
\vspace{0.8cm}
{\large 11 October 2026\par}
\vspace{0.8cm}
\begin{abstract}
We determine, at every weight, the complete projection of the inverse-Gamma
Bell hierarchy onto functional monomials involving digamma and its first
three derivatives. A single entire integral formula gives all coefficients
of $u_0^a u_1^k u_2^b u_3^n$ for $n\ge2$, with explicit parity zeros and a
finite two-binomial rational expression. The square case resolves the
coupled next-Gamma-sector question in the inspected ProveIt continuation;
the cubic and higher powers of $u_3$ go beyond that question. The proof uses
the Euler variational derivative, an injective triangular operator, and
four explicit polynomial identities, rather than extrapolation from finite
weight calculations. A beta-density formula also replaces the preceding
pure-$u_2$ recurrence by a nonrecursive integral.

The same endpoint-vanishing polynomials generate finite Hurwitz-zeta and
Lerch formulas, all-order generalized Stieltjes averages and anchored
primitives, and ordinary polylogarithm sums with harmonic endpoint values.
Several weighted digamma and first-Stieltjes integrals are evaluated
explicitly. Canonical total-derivative primitives are retained when the
formal identities are transported to Gamma and harmonic numbers.

Independent exact reductions check 973 mixed coefficients through weight
30; symbolic checks verify the four all-index density identities. Thirty
high-precision numerical diagnostics are supplied separately. An
arbitrary-derivative-order extension is formulated as an unproved
conjecture, with its computational evidence clearly separated. Global
priority and exhaustive audit of the incoming archives are not claimed.
\end{abstract}
\vfill
\noindent\textbf{Included artifacts.} Complete and modular \LaTeX\ sources,
compiled PDF, rational verification programs and certificates, numerical
diagnostics, source provenance, and repository integration notes.
\end{titlepage}
\hypersetup{pageanchor=true}
\pagenumbering{roman}
\tableofcontents
\clearpage\pagenumbering{arabic}
% BEGIN INCLUDED FILE: sections/01_scope.tex
\section{The research question and the scope of the answer}
\label{sec:scope}

The inspected ProveIt manuscript develops exact identities among
polylogarithms, Gamma derivatives, harmonic tails, and Hurwitz--Stieltjes
coordinates \cite{RepoMain}. Its differential-polynomial continuation
already proves a canonical integration-by-parts normal form and a complete
recurrence for the variables $u_0,u_1,u_2$ \cite{RepoQuadratic}. A later
continuation treats square sectors on an affine background and explicitly
asks for every coefficient
\begin{equation}\label{eq:source-question}
 [u_0^a u_1^k u_2^b u_3^2]N_{a+2k+3b+8},\qquad b>0,
\end{equation}
with parity and the endpoint primitive retained \cite[Question 1]{RepoQuestions}.
Here $u_r$ represents the $r$th derivative of digamma, and $N_m$ is the
canonical normal representative of
$(-1)^m\Gamma(x)(\mathrm d/\mathrm dx)^m\Gamma(x)^{-1}$.

We answer \eqref{eq:source-question} without restricting $a,k,b$. More
strongly, we allow any number $n\ge2$ of $u_3$ factors. The answer is a
finite rational sum whose summands have four consecutive odd factors in
the denominator. Its underlying integral is a beta moment with the
endpoint factor $(1-v^2)^3$. This factor is not decorative: its three
vanishing endpoint derivatives are exactly what makes the third-order
variational transport close.

\subsection{A preview in ordinary coefficients}
For $m=a+2k+3b+4n$ and $n\ge2$, the coefficient vanishes when $k+n$ is
odd. When $k+n$ is even it is
\begin{equation}\label{eq:preview}
\begin{split}
 &[u_0^a u_1^k u_2^b u_3^n]N_m\\
 &\quad=-\frac{m!}{192\,a!\,2^k k!\,b!\,n!\,24^b48^{n-2}}
 \sum_{j=0}^{b}\sum_{\ell=0}^{n-2}
 \frac{\binom bj3^j\binom{n-2}{\ell}}
 {\displaystyle\prod_{r=0}^{3}(k+n+2j+2\ell-1+2r)}.
\end{split}
\end{equation}
For example, the coefficients of $u_2u_3^2$, $u_2^2u_3^2$, and
$u_1u_3^3$ are respectively $-55$, $-3640$, and $-3185/3$. These are
coefficients in the specified rational normal basis, not assertions of
arithmetic independence after special-function substitution.

\subsection{What is inherited, and what is established here}
The normal-form decomposition and the finite-jet Euler principle are
inherited. We reproduce their elementary proofs so that the article does
not depend on trusting a computation or an inaccessible manuscript.
The underlying functional-monomial framework is credited in the source to
Guo, Regensburger, and Rosenkranz \cite{GRR}. The preceding pure-$u_2$
recurrence is also inherited; its beta-density solution and the complete
cubic-background formula are the developments proved below.

The polynomial-weight Hurwitz and Lerch integration-by-parts mechanism
uses classical special-function identities \cite{DLMF25,DLMF14}. We do not
claim that mechanism, the standard Gamma--zeta relations, or elementary
partial fractions as new. Their application to the newly determined mixed
coefficient kernels gives the exact families and concrete evaluations in
Sections~\ref{sec:hurwitz} and \ref{sec:polylog}. Historical priority for
these specializations has not been exhaustively investigated.

The incoming directory was inventoried during source selection, but its
binary packages were not exhaustively read or replayed. In particular,
the conclusions here are not a claim that every current delivery has been
reconciled. The provenance file distinguishes the inspected text from
archive names seen in the directory. Since the repository changed during
retrieval, we record file-level identifiers where available instead of
claiming a single atomic repository snapshot.

\subsection{Conventions}
All formal polynomial identities are over $\Q$. We write
\[
 \avg{f(v)}=\frac12\int_{-1}^1f(v)\dd v,
 \qquad [t^m]F(t)\text{ for ordinary coefficient extraction}.
\]
The exponential generating functions carry their displayed factorials.
The notation $\zeta^{(r)}(s,a)$ always means an \emph{order} derivative,
$\partial_s^r\zeta(s,a)$; $\psi^{(r)}(a)$ always means an \emph{argument}
derivative. Stieltjes constants use the convention
\begin{equation}\label{eq:stieltjes-convention}
 \zeta(1+\epsilon,a)=\frac1\epsilon+
 \sum_{q\ge0}\frac{(-1)^q}{q!}\gamma_q(a)\epsilon^q,
 \qquad \gamma_0(a)=-\psi(a).
\end{equation}
For analytic integral formulas we initially impose $\Re a>1$, so that
$a+v$ stays in the right half-plane. All polylogarithm series are initially
used on $|z|<1$, eliminating branch ambiguity. Extensions and removable
singularities are stated where used.

% END INCLUDED FILE: sections/01_scope.tex

% BEGIN INCLUDED FILE: sections/02_normal_forms.tex
\section{Canonical normal forms and the Euler principle}
\label{sec:normal}

Let
\[
 \Rcal=\Q[u_0,u_1,u_2,\ldots],\qquad
 \D u_j=u_{j+1},\qquad \wt(u_j)=j+1.
\]
The factor degree of a monomial is its ordinary total degree.
A monomial is \emph{functional} if it is a power of $u_0$, including $1$,
or if its largest positive derivative index occurs at least twice.
Denote the rational span of these monomials by $\mathscr F$.

\begin{proposition}[Canonical decomposition]\label{prop:normal}
Every $P\in\Rcal$ has a unique decomposition
\begin{equation}\label{eq:decomp}
 P=\NF(P)+\D\Prim(P),\qquad
 \NF(P)\in\mathscr F,\qquad [1]\Prim(P)=0.
\end{equation}
Both maps are rational and linear. They preserve factor degree; on a
homogeneous input of weight $m$, $\NF(P)$ has weight $m$ and $\Prim(P)$
has weight $m-1$.
\end{proposition}
\begin{proof}
Suppose that a monomial $M$ has a unique largest derivative $u_r$,
$r\ge1$. Replace that factor by $u_{r-1}$, obtaining $T$, and let $c$
be the multiplicity of $u_{r-1}$ in $T$. Then
\[
 M=\frac1c\D T+\left(M-\frac1c\D T\right).
\]
The remaining monomials have strictly smaller largest derivative index.
Repeated application terminates and constructs both maps. The asserted
gradings are preserved at every step.

For uniqueness, let a nonconstant $H$ have largest index $r$. In $\D H$
the term $u_{r+1}\partial_{u_r}H$ is nonzero and consists entirely of
monomials with unique largest index $r+1$. It cannot be canceled by the
other terms of $\D H$, which have largest index at most $r$, or by a
functional polynomial. Thus $\D H\notin\mathscr F$. This proves
uniqueness once the constant term of $H$ is fixed.
\end{proof}

The decomposition is a direct sum of vector spaces, not of algebras.
In particular, $\NF$ is not multiplicative: $\NF(u_1)=0$ while
$\NF(u_1^2)=u_1^2$. The constructive proof gives a complete endpoint
primitive, and the supplied rational program implements that proof.

\subsection{The inverse-Gamma Bell hierarchy}
Define
\begin{align}
 A_0&=1,& A_{m+1}&=u_0A_m-\D A_m,\label{eq:Bell-recur}\\
 \Acal(t)&=\sum_{m\ge0}A_m\frac{t^m}{m!}
 =\exp\left(\sum_{r\ge0}\frac{(-1)^ru_rt^{r+1}}{(r+1)!}\right),
 \label{eq:Bell}\end{align}
and put $N_m=\NF(A_m)$ and $H_m=\Prim(A_m)$.
Formula~\eqref{eq:Bell} follows either by differentiating the exponential
or by the formal Taylor formula. Under $u_r=\psi^{(r)}(x)$,
\begin{equation}\label{eq:Gamma-realization}
 A_m(x)=(-1)^m\Gamma(x)\frac{\mathrm d^m}{\mathrm dx^m}\frac1{\Gamma(x)},
 \qquad \Acal(t,x)=\frac{\Gamma(x)}{\Gamma(x-t)}.
\end{equation}
These identifications use $\psi=\Gamma'/\Gamma$ and its derivatives
\cite{DLMF5}. For every fixed $m$ they are ordinary analytic identities
away from Gamma poles.

Translation $u_0\mapsto u_0+c$ commutes with $\D$ and preserves
functional monomials. By uniqueness it commutes with $\NF$. Consequently
$\NF\Acal(t)$ has an Appell factor $e^{u_0t}$. A second useful identity is
\begin{equation}\label{eq:centered}
 \NF\Acal(t)=\NF\exp\left(
 \sum_{j\ge0}\frac{u_{2j}t^{2j+1}}{2^{2j}(2j+1)!}\right).
\end{equation}
Indeed the exponential before normalization on the right is
$e^{t\D/2}\Acal(t)$. Every coefficient of
$(e^{t\D/2}-1)\Acal(t)$ is a total derivative. The exponent in
\eqref{eq:centered} has odd weights only, so every retained monomial has
factor degree congruent to its weight modulo two. For
$u_0^a u_1^k u_2^b u_3^n$ this says that $k+n$ must be even.

\subsection{Euler derivatives and finite jets}
Define
\[
 \Euler(P)=\sum_{r\ge0}(-\D)^r\partial_{u_r}P.
\]
Every sum is finite on a polynomial. The commutator relation
$\partial_{u_r}\D=\D\partial_{u_r}+\partial_{u_{r-1}}$, with the
last term zero at $r=0$, gives $\Euler\D=0$ by cancellation.

Let $\sigma_d$ set all $u_r$ with $r>d$ to zero and define
$\delta_d=\sum_{r=0}^{d-1}u_{r+1}\partial_{u_r}$.

\begin{lemma}[Finite-jet variational identity]\label{lem:finitejet}
For every $P\in\Rcal$,
\begin{equation}\label{eq:finitejet}
 \sigma_d\Euler(P)=
 \sum_{r=0}^d(-\delta_d)^r\partial_{u_r}\bigl(\sigma_d\NF(P)\bigr).
\end{equation}
Furthermore,
\begin{equation}\label{eq:Euler-A}
 \Euler\Acal(t)=\int_0^t e^{s\D}\Acal(t)\dd s.
\end{equation}
The integral in \eqref{eq:Euler-A} is coefficientwise formal.
\end{lemma}
\begin{proof}
Since $\Euler\D=0$, replace $P$ by $\NF(P)$ on the left of
\eqref{eq:finitejet}. A normal monomial removed by $\sigma_d$ has at
least two factors of its highest index, which exceeds $d$. One partial
derivative removes at most one such factor; subsequent applications of
$\D$ preserve a factor of index greater than $d$. Hence its specialization
vanishes. For retained monomials, the ideal generated by
$u_{d+1},u_{d+2},\ldots$ is $\D$-stable, and specialization changes
$\D$ to $\delta_d$. This proves \eqref{eq:finitejet}.

By \eqref{eq:Bell},
$\partial_{u_r}\Acal(t)=(-1)^rt^{r+1}\Acal(t)/(r+1)!$. Therefore
\[
 \Euler\Acal(t)=\sum_{r\ge0}\frac{t^{r+1}}{(r+1)!}\D^r\Acal(t),
\]
which is the formal integral in \eqref{eq:Euler-A}.
\end{proof}

Normal reduction must precede truncation. For example,
$\sigma_1(u_0u_2)=0$, whereas
$\sigma_1\NF(u_0u_2)=-u_1^2$.
The finite-jet lemma, not commutation of these operations, is the reason
one can obtain a closed low-derivative equation.

% END INCLUDED FILE: sections/02_normal_forms.tex

% BEGIN INCLUDED FILE: sections/03_quadratic.tex
\section{A beta-density solution of the quadratic background}
\label{sec:quadratic}

We use the dimensionless variables and polynomials
\begin{equation}\label{eq:basic-variables}
 y=\frac{u_1t^2}{2},\quad z=u_2t^3,\quad w=u_3t^4,\qquad
 h(v)=\frac{1+3v^2}{24},\quad k(v)=\frac{v+v^3}{48},
\end{equation}
with $P_r(v)=(1-v^2)^r$. In particular $h'=v/4$ and $k'=h/2$.
Write
\[
 \sum_{m\ge0}\sigma_2N_m\frac{t^m}{m!}
 =e^{u_0t}F_0(y,z),\qquad F_0(y,z)=\sum_{b\ge0}z^bf_b(y).
\]
The previous continuation gives a differential equation and rational
recurrence for these functions \cite{RepoQuadratic}. We solve that
recurrence by a single nonrecursive beta density.

\begin{theorem}[Pure-$u_2$ density]\label{thm:quadratic}
Let
\begin{equation}\label{eq:faff}
 f_{\mathrm{aff}}(y)=\frac12\left(
 \cosh y+\frac{\sinh y}{y}-y\Shi(y)\right),
 \qquad \Shi(y)=\int_0^y\frac{\sinh v}{v}\dd v.
\end{equation}
All singularities displayed at zero are removable. Then $f_0=f_{\mathrm{aff}}$,
$f_1=0$, and, for every $b\ge2$,
\begin{equation}\label{eq:fb-density}
 \boxed{\displaystyle f_b(y)=\frac1{192b!}
  \avg{P_2(v)h(v)^{b-2}e^{yv}}.}
\end{equation}
Equivalently, with
$\Erem(x)=(e^x-1-x)/x^2$ and $\Erem(0)=1/2$,
\begin{equation}\label{eq:F0-compact}
 F_0(y,z)=f_{\mathrm{aff}}(y)+\frac{z^2}{192}
 \avg{P_2(v)e^{yv}\Erem(zh(v))}.
\end{equation}
This is entire jointly in $y,z$, and
\begin{equation}\label{eq:F0zz}
 \partial_z^2F_0(y,z)=\frac1{192}\avg{P_2(v)e^{yv+zh(v)}}.
\end{equation}
\end{theorem}

\begin{proof}
Let $A=2y+(z/2)\partial_y$. Lemma~\ref{lem:finitejet}, specialized to
the quadratic background, gives
\begin{equation}\label{eq:quadratic-PDE}
 F_0-\frac12A\partial_yF_0+A^2\partial_zF_0
 =\avg{e^{yv+zh(v)}}.
\end{equation}
To see the right-hand side, set $s=t(1+v)/2$ in
\eqref{eq:Euler-A}; the exponent after specialization is
$u_0t+yv+zh(v)$. The two sides of \eqref{eq:quadratic-PDE} are the
Euler identity divided by $te^{u_0t}$.

Since a single $u_2$ is not functional, $f_1=0$. Coefficient $z^0$
in \eqref{eq:quadratic-PDE} yields
$f_0-yf_0'=\sinh y/y$. The unique solution with $f_0(0)=1$ and no
linear term is \eqref{eq:faff}. In particular,
\begin{equation}\label{eq:faff-second}
 f_{\mathrm{aff}}''(y)=-\frac12\avg{(1-v^2)e^{yv}}.
\end{equation}
For $n\ge1$, the remaining coefficients in the PDE give
\begin{equation}\label{eq:quadratic-rec}
\begin{split}
 4(n+1)y^2f_{n+1}
 =\frac1{n!}\avg{h^ne^{yv}}
 -(n+1)f_n-(2n-1)yf_n'
 -\frac{n-2}{4}f_{n-1}''.
\end{split}
\end{equation}
The case $n=1$, using \eqref{eq:faff-second}, gives
$f_2=\avg{P_2e^{yv}}/384$. For instance, this follows immediately
from $P_2''=4(3v^2-1)$ by integrating twice by parts. The endpoint
values of $P_2$ and $P_2'$ are zero.

For $n\ge2$ the elementary polynomial identity
\begin{equation}\label{eq:pure-density-identity}
\begin{split}
 4(P_2h^{n-1})''={}&192h^n-(n+1)P_2h^{n-2}
 +(2n-1)(vP_2h^{n-2})'\\
 &-\frac{n(n-2)}4v^2P_2h^{n-3}
\end{split}
\end{equation}
proves that \eqref{eq:fb-density} satisfies \eqref{eq:quadratic-rec}.
At $n=2$ the last term is understood as zero. For all $n\ge3$, divide
by $h^{n-3}$, differentiate with $h'=v/4$, and substitute
$24h=1+3v^2$; both sides become the same polynomial in $n$ and $v$.
This is a direct all-index algebraic identity, not an inference from
checking finitely many $n$.

To pass from \eqref{eq:pure-density-identity} to
\eqref{eq:quadratic-rec}, integrate against $e^{yv}$ and transfer
one or two $v$ derivatives. Endpoint terms vanish because of $P_2$.
The terms involving $v$ and $v^2$ give $f_n'$ and $f_{n-1}''$.
After multiplication by $192n!$, the result is exactly the recurrence.
Multiplication by $y^2$ is injective on formal power series, so the
initial values determine the sequence uniquely. This proves the density
formula. Summing $z^b$ gives \eqref{eq:F0-compact}, and differentiating
twice gives \eqref{eq:F0zz}. Compact integration of entire functions
proves joint entire convergence.
\end{proof}

For even $k$, extracting the coefficient of $y^k$ gives a particularly
simple rational replacement for the source recurrence:
\begin{equation}\label{eq:quadratic-beta}
 k![y^k]f_b(y)=\frac1{24^{b-1}b!}
 \sum_{j=0}^{b-2}
 \frac{\binom{b-2}{j}3^j}{(k+2j+1)(k+2j+3)(k+2j+5)}.
\end{equation}
For odd $k$ it is zero. Formula~\eqref{eq:quadratic-beta} follows by
expanding $(1+3v^2)^{b-2}$ and using
$\int_0^1v^K(1-v^2)^2\dd v=8/((K+1)(K+3)(K+5))$.

% END INCLUDED FILE: sections/03_quadratic.tex

% BEGIN INCLUDED FILE: sections/04_cubic.tex
\clearpage
\section{The complete cubic-background formula}
\label{sec:cubic}

We now determine the entire projection with highest retained derivative
index three. The words ``cubic background'' refer to the fact that a
formal shift of $u_0$ is a cubic polynomial when $u_r=0$ for $r>3$.
They do not restrict the power of $u_3$ in the answer.

\begin{theorem}[Complete cubic normal form]\label{thm:main}
With the notation of \eqref{eq:basic-variables},
\begin{equation}\label{eq:main-generating}
 \sum_{m\ge0}\sigma_3N_m\frac{t^m}{m!}=e^{u_0t}F(y,z,w),
\end{equation}
where
\begin{equation}\label{eq:main}
\boxed{\begin{aligned}
 F(y,z,w)={}&f_{\mathrm{aff}}(y)
 +\frac{z^2}{192}\avg{P_2e^{yv}\Erem(zh)}\\
 &-\frac{w^2}{9216}\avg{P_3e^{yv+zh}\Erem(wk)}.
\end{aligned}}
\end{equation}
This function is entire jointly in $y,z,w$. Writing
$F=\sum_{n\ge0}w^nF_n(y,z)$, its coefficients are $F_0$ from
Theorem~\ref{thm:quadratic}, $F_1=0$, and
\begin{equation}\label{eq:Fn-density}
 F_n(y,z)=-\frac1{9216n!}
 \avg{P_3(v)k(v)^{n-2}e^{yv+zh(v)}},\qquad n\ge2.
\end{equation}
\end{theorem}

\subsection{The closed Euler equation}
Set
\[
 A=2y+\frac z2\partial_y,\qquad B=\partial_z,\qquad
 C=\frac12\partial_y,\qquad L=A+wB.
\]
On a function $e^{u_0t}G(y,z,w)$ the derivation $\delta_3$ acts as
$e^{u_0t}t^{-1}LG$. Lemma~\ref{lem:finitejet} therefore gives
\begin{equation}\label{eq:cubic-PDE}
 F-\frac12L F_y+L^2F_z-L^3F_w
 =\avg{e^{yv+zh(v)+wk(v)}}.
\end{equation}
In detail, after $s=t(1+v)/2$ in \eqref{eq:Euler-A}, the exponent of
the shifted Bell exponential is
\[
 u_0t+\frac{u_1t^2v}{2}
 +\frac{u_2t^3(1+3v^2)}{24}
 +\frac{u_3t^4(v+v^3)}{48}.
\]
This establishes the right-hand side of \eqref{eq:cubic-PDE} with all
normalizations fixed.

The commutators $[B,A]=C$ and $[C,A]=1$ give
\begin{align*}
 L^2&=A^2+w(2AB+C)+w^2B^2,\\
 L^3&=A^3+w(3A^2B+3AC+1)
          +w^2(3AB^2+3CB)+w^3B^3.
\end{align*}
Put
\[
 R_n(y,z)=\frac1{n!}\avg{k(v)^ne^{yv+zh(v)}}.
\]
Comparison of $w^n$ in \eqref{eq:cubic-PDE} gives, for $n\ge1$,
\begin{equation}\label{eq:cubic-rec}
\begin{split}
 (n+1)A^3F_{n+1}={}&(1-n)F_n-\frac{1+3n}{2}A\partial_yF_n
 +(1-3n)A^2\partial_zF_n\\
 &+(5-3n)A\partial_z^2F_{n-1}
 +\frac{3(1-n)}2\partial_y\partial_zF_{n-1}\\
 &+(3-n)\partial_z^3F_{n-2}-R_n.
\end{split}
\end{equation}
Negative subscripts mean zero. The $w^0$ equation is exactly the
quadratic equation because the normal basis excludes a single $u_3$.

\begin{lemma}[Uniqueness of the cubic hierarchy]\label{lem:injective}
The operator $A^3$ is injective on $\C[[y,z]]$. Consequently
\eqref{eq:cubic-rec}, together with the specified $F_0$ and $F_1=0$,
determines every $F_n$ uniquely.
\end{lemma}
\begin{proof}
If $G$ is nonzero, let $j$ be its least nonzero $z$-degree and write that
coefficient as $g(y)$. The coefficient of $z^j$ in $A^3G$ is
$8y^3g(y)$, which is nonzero. Indeed every use of $(z/2)\partial_y$
instead of $2y$ raises $z$-degree. This proves injectivity and then the
successive uniqueness assertion.
\end{proof}

\subsection{Why the beta density solves every recurrence}
Let $E(v)=e^{yv+zh(v)}$. The key transport relation is
\begin{equation}\label{eq:transport}
 A E(v)=2\partial_v E(v),\qquad
 A^j\avg{Q(v)E(v)}=(-2)^j\avg{Q^{(j)}(v)E(v)}.
\end{equation}
The second equality holds when $Q,Q',\ldots,Q^{(j-1)}$ vanish at
both endpoints. For $Q$ divisible by $P_3$, it applies for $j\le3$.
The coefficients of $A$ are independent of $v$, so iteration of the
first equality and ordinary integration by parts prove the second.

\begin{proof}[Proof of Theorem~\ref{thm:main}]
We verify that \eqref{eq:Fn-density} satisfies
\eqref{eq:cubic-rec}. Uniqueness then identifies it with the actual
normal projection, so no unproved assumption about a formal solution is
needed.

For $n=1$, the recurrence reduces to
\[
 2A^3F_2=2A\partial_z^2F_0-R_1.
\]
Using \eqref{eq:F0zz} and \eqref{eq:transport}, this is the polynomial
identity
\begin{equation}\label{eq:cubic-base1}
 \frac{P_3'''}{1152}=-\frac{P_2'}{48}-k.
\end{equation}
For $n=2$, the required identity is
\begin{equation}\label{eq:cubic-base2}
 -8(P_3k)'''=-P_3+7(vP_3)'-20(hP_3)''-96P_2h+9216k^2.
\end{equation}
Both follow by differentiating the displayed polynomials.

For every $n\ge3$, substituting \eqref{eq:Fn-density} in the
recurrence, multiplying by $-9216n!$, and using
\eqref{eq:transport} reduces the assertion to
\begin{equation}\label{eq:cubic-density-identity}
\begin{split}
 -8(P_3k^{n-1})'''={}&(1-n)P_3k^{n-2}
 +(1+3n)(vP_3k^{n-2})'\\
 &+4(1-3n)(hP_3k^{n-2})''
 -2n(5-3n)(h^2P_3k^{n-3})'\\
 &+\frac32n(1-n)vhP_3k^{n-3}
 +n(n-1)(3-n)h^3P_3k^{n-4}\!\\
 &+9216k^n.
\end{split}
\end{equation}
At $n=3$ the term with factor $3-n$ is zero; no negative power is
part of the asserted identity. For $n\ge4$ all terms are polynomials.

Here is a finite algebraic way to verify \eqref{eq:cubic-density-identity}
for an indeterminate $n$, making the all-index step explicit. Define
\[
 \mathcal G_{a,q}(U)=U'+a\frac{q'}qU.
\]
On the open set where $k\ne0$, divide \eqref{eq:cubic-density-identity}
by $k^{n-4}$ and replace its left side by
$-8\mathcal G_{n-4,k}^3(P_3k^3)$. The terms on its right become,
in order,
\begin{align*}
 &(1-n)P_3k^2,
 &&(1+3n)\mathcal G_{n-4,k}(vP_3k^2),\\
 &4(1-3n)\mathcal G_{n-4,k}^2(hP_3k^2),
 &&-2n(5-3n)\mathcal G_{n-4,k}(h^2P_3k),\\
 &\tfrac32n(1-n)vhP_3k,
 &&n(n-1)(3-n)h^3P_3,
 \qquad 9216k^4.
\end{align*}
Substitute $P_3=(1-v^2)^3$, $24h=1+3v^2$, and $48k=v+v^3$.
Expanding the derivatives and collecting powers of $n$ makes the
difference identically zero in $\Q(v)[n]$. Only three derivatives are
needed, independently of $n$. Clearing denominators proves the identity
on the entire polynomial domain. The same calculation at $n=3$, or its
polynomial continuation, proves that case with the stated zero-term
convention. The supplied script records these exact symbolic operations.

All integrations by parts in this verification have zero endpoint terms:
the worst derivative transfer is of order three, and its density has the
factor $P_3$. The exceptional terms involving $F_0$ occur only at
$n=1,2$ and were handled separately. Thus the recurrence is verified at
every index. Lemma~\ref{lem:injective} proves
\eqref{eq:Fn-density}. Finally,
$\sum_{n\ge2}w^nk^{n-2}/n!=w^2\Erem(wk)$ gives
\eqref{eq:main}. The integrands are entire and uniformly bounded on
compact parameter sets and the compact interval, proving joint entirety.
\end{proof}

\begin{remark}[Proof status of the algebraic step]
The four identities \eqref{eq:pure-density-identity},
\eqref{eq:cubic-base1}, \eqref{eq:cubic-base2}, and
\eqref{eq:cubic-density-identity} are explicit finite polynomial or
rational-function identities. Their symbolic verification treats $n$ as
an indeterminate. The separate checks at weights at most 30 are useful
regression tests but are not the reason that the theorem holds for every
weight.
\end{remark}

% END INCLUDED FILE: sections/04_cubic.tex

% BEGIN INCLUDED FILE: sections/05_coefficients.tex
\section{Finite coefficient formulas and explicit mixed families}
\label{sec:coefficients}

\begin{corollary}[All mixed $u_3$ coefficients]\label{cor:coefficients}
Let $a,k,b\ge0$, $n\ge2$, and $m=a+2k+3b+4n$. If $k+n$ is odd,
then $[u_0^au_1^ku_2^bu_3^n]N_m=0$. If $k+n$ is even, then
\begin{equation}\label{eq:coefficient-integral}
\begin{split}
 &[u_0^au_1^ku_2^bu_3^n]N_m\\
 &=-\frac{m!}{9216\,a!2^kk!b!n!24^b48^{n-2}}
 \int_0^1v^{k+n-2}(1-v^2)^3
       (1+3v^2)^b(1+v^2)^{n-2}\dd v.
\end{split}
\end{equation}
Equivalently, the finite rational formula \eqref{eq:preview} holds.
Every coefficient not forced to vanish by parity is strictly negative.
\end{corollary}
\begin{proof}
Expand $e^{u_0t}$, $e^{yv}$, $e^{zh(v)}$, and $\Erem(wk(v))$ in
\eqref{eq:main}. The conversion from $y^k$ to $u_1^k$ contributes
$2^{-k}$, and conversion to $N_m$ contributes $m!$. The remaining
integrand is odd precisely when $k+n$ is odd; otherwise its average on
$[-1,1]$ is its integral on $[0,1]$. This proves
\eqref{eq:coefficient-integral}. Expand the two binomials and use
\begin{equation}\label{eq:beta-four}
 \int_0^1v^K(1-v^2)^3\dd v
 =\frac{48}{(K+1)(K+3)(K+5)(K+7)},\qquad K>-1.
\end{equation}
Formula~\eqref{eq:beta-four} follows by $x=v^2$ and three elementary
integrations of the beta integral, or by expanding $(1-v^2)^3$.
Here $K=k+n-2+2j+2\ell\ge0$, so every denominator is positive.
This proves the finite formula and the sign assertion.
\end{proof}

\subsection{The source square-sector question}
For $n=2$ and even $k$, \eqref{eq:preview} becomes
\begin{equation}\label{eq:square-answer}
\boxed{\begin{aligned}
 &[u_0^au_1^ku_2^bu_3^2]N_{a+2k+3b+8}\\
 &=-\frac{(a+2k+3b+8)!}{384\,a!2^kk!b!24^b}
 \sum_{j=0}^b
 \frac{\binom bj3^j}
 {(k+2j+1)(k+2j+3)(k+2j+5)(k+2j+7)}.
\end{aligned}}
\end{equation}
It is zero for odd $k$. Thus every $b>0$ in the specific question
\eqref{eq:source-question} is covered; no weight-dependent linear
system remains to be solved.

Define the normalized square coefficient for even $k$ by
\[
 g_b(k)=k![y^kz^bw^2]F(y,z,w).
\]
Then
\begin{equation}\label{eq:gb-density}
 g_b(k)=-\frac1{18432b!24^b}
 \int_0^1v^k(1-v^2)^3(1+3v^2)^b\dd v.
\end{equation}
For example,
\begin{align}
 g_0(k)&=-\frac1{384(k+1)(k+3)(k+5)(k+7)},\label{eq:g0}\\
 g_1(k)&=-\frac1{2304(k+1)(k+5)(k+7)(k+9)},\label{eq:g1}\\
 g_2(k)&=-\frac{(k+2)(k+6)}{27648\prod_{j=0}^5(k+2j+1)},\label{eq:g2}\\
 g_3(k)&=-\frac{k^3+15k^2+62k+63}{497664\prod_{j=0}^6(k+2j+1)}.
 \label{eq:g3}
\end{align}
These rational functions are asserted as coefficient formulas at even
nonnegative integers $k$. Their analytic continuation in $k$ is also
provided by the beta integral for $\Re k>-1$, but the coefficient at an
odd power of $y$ remains zero.

\subsection{Nonlinear sectors beyond the square}
The complete formula also gives genuinely coupled cubic and quartic
sectors. A small sample, with no $u_0$ factor, is
\begin{center}
\begin{tabular}{rll}
\toprule
Weight & Functional monomial & Coefficient in $N_m$\\
\midrule
8 & $u_3^2$ & $-1$\\
11 & $u_2u_3^2$ & $-55$\\
12 & $u_1^2u_3^2$ & $-165$\\
14 & $u_2^2u_3^2$ & $-3640$\\
14 & $u_1u_3^3$ & $-3185/3$\\
16 & $u_3^4$ & $-10325/3$\\
17 & $u_1u_2u_3^3$ & $-1023400/3$\\
19 & $u_2u_3^4$ & $-1639225$\\
\bottomrule
\end{tabular}
\end{center}
All coefficients with an additional power $u_0^a$ follow immediately
from the Appell factorial ratio. The rational denominators in this
basis need not be integral, even though the Bell polynomial before
normal reduction has integral coefficients.

The formulas also reveal a useful structural distinction. The pure
highest-index-two sector has positive beta density, whereas the entire
highest-index-three sector has negative beta density, subject to parity.
This observation motivates the all-order conjecture in
Section~\ref{sec:conjecture}; it does not by itself establish that
conjecture.

% END INCLUDED FILE: sections/05_coefficients.tex

% BEGIN INCLUDED FILE: sections/06_gamma_harmonic.tex
\section{Gamma identities, harmonic values, and retained primitives}
\label{sec:gamma}

The normal-form theorem is an identity in a free differential algebra.
Its analytic meaning must keep both the primitive in
\eqref{eq:decomp} and the distinction between a projection and the full
normal polynomial.

\begin{theorem}[Transport with the exact endpoint]\label{thm:endpoint}
Let $0<A<B$, $m\ge0$, and write
$U(x)=(\psi(x),\psi'(x),\ldots)$. Then
\begin{equation}\label{eq:Gamma-endpoint}
\begin{split}
 \int_A^B N_m(U(x))\dd x
 ={}&(-1)^m\int_A^B\Gamma(x)
       \frac{\mathrm d^m}{\mathrm dx^m}\frac1{\Gamma(x)}\dd x\\
 &-H_m(U(B))+H_m(U(A)).
\end{split}
\end{equation}
The polynomial $H_m$ is computed by the terminating rational construction
in Proposition~\ref{prop:normal}. In particular, the formula gives an
anchored primitive for every normal combination $N_m$ without suppressing
an integration-by-parts boundary constant.
\end{theorem}
\begin{proof}
Substitute $u_r=\psi^{(r)}(x)$ into
$A_m=N_m+\D H_m$. Under this substitution $\D$ is ordinary argument
differentiation. Use \eqref{eq:Gamma-realization} and integrate over the
compact interval. All functions are analytic there.
\end{proof}

For use with the compact cubic formula, let
$R_m=N_m-\sigma_3N_m$. The precise projected identity is therefore
\begin{equation}\label{eq:projected-endpoint}
\begin{split}
 \int_A^B(\sigma_3N_m)(U(x))\dd x
 ={}&(-1)^m\int_A^B\Gamma(x)(\Gamma^{-1})^{(m)}(x)\dd x\\
 &-[H_m(U(x))]_A^B-\int_A^BR_m(U(x))\dd x.
\end{split}
\end{equation}
One must not erase the last term when the higher-derivative sector is
nonzero. For $m\le9$ it vanishes by weight, but it need not vanish from
weight ten onward.

\subsection{Two explicit normal polynomials}
The weight-eight identity is
\begin{equation}\label{eq:N8}
 N_8=u_0^8-70u_0^4u_1^2+28u_0^2u_2^2-7u_1^4-u_3^2.
\end{equation}
For this case the complete endpoint polynomial is
\begin{equation}\label{eq:H8}
\begin{split}
 H_8={}&-4u_0^7+56u_0^5u_1-70u_0^4u_2-140u_0^3u_1^2
       +56u_0^3u_3\\
 &+252u_0^2u_1u_2-28u_0^2u_4+112u_0u_1^3
       -112u_0u_1u_3\\
 &-84u_0u_2^2+8u_0u_5-98u_1^2u_2+20u_1u_4+36u_2u_3-u_6.
\end{split}
\end{equation}
Direct differentiation verifies $A_8=N_8+\D H_8$.
Equations~\eqref{eq:Gamma-endpoint}, \eqref{eq:N8}, and \eqref{eq:H8}
are an ordinary finite-interval integral identity involving digamma,
trigamma, and higher polygammas, with an entirely explicit endpoint.

At weight eleven,
\begin{equation}\label{eq:N11}
\begin{split}
 N_{11}={}&u_0^{11}-330u_0^7u_1^2+462u_0^5u_2^2
 -1155u_0^3u_1^4-165u_0^3u_3^2\\
 &+550u_0^2u_2^3+990u_0u_1^2u_2^2+11u_0u_4^2-55u_2u_3^2.
\end{split}
\end{equation}
Thus $R_{11}=11u_0u_4^2$. This gives a concrete test of why the
projection alone cannot be substituted for the full $N_{11}$ in an
unqualified Gamma integral. The coefficient $-55$ is the first coupled
case of \eqref{eq:square-answer}. The full $H_{11}$ and other exact
primitives are supplied as rational JSON certificates.

\subsection{Generalized harmonic numbers}
Let $H_N^{(r)}=\sum_{j=1}^Nj^{-r}$, $H_N=H_N^{(1)}$, and let
$\gamma$ denote Euler's constant. The standard shift formulas give
\begin{equation}\label{eq:harmonic-substitution}
 u_0=H_N-\gamma,\qquad
 u_r=(-1)^{r+1}r!\bigl(\zeta(r+1)-H_N^{(r+1)}\bigr),\quad r\ge1,
\end{equation}
at $x=N+1$ \cite{DLMF5,DLMF25}. Let $U_N$ denote this tuple. Then
\begin{equation}\label{eq:harmonic-identity}
 N_m(U_N)+(\D H_m)(U_N)
 =m![t^m]\frac{\Gamma(N+1)}{\Gamma(N+1-t)}.
\end{equation}
The right side also equals
\begin{equation}\label{eq:harmonic-generator}
 m![t^m]\exp\left((H_N-\gamma)t-
 \sum_{r\ge2}\frac{\zeta(r)-H_N^{(r)}}r\,t^r\right).
\end{equation}
This follows by substituting \eqref{eq:harmonic-substitution} into
\eqref{eq:Bell}. It is initially a local analytic series in $t$ and an
identity of every finite coefficient. The derivative $\D H_m$ is formed
in the free differential algebra \emph{before} substitution; it is not a
finite difference in $N$.

Consequently \eqref{eq:square-answer} and the nonlinear coefficients
specify exact mixed harmonic-tail contributions to every identity
\eqref{eq:harmonic-identity}. This is a systematic family of identities,
not a claim that the individual harmonic products have been shown
arithmetically independent or that every integral of one isolated
polygamma monomial has been separately evaluated.

% END INCLUDED FILE: sections/06_gamma_harmonic.tex

% BEGIN INCLUDED FILE: sections/07_hurwitz_stieltjes.tex
\section{Hurwitz endpoints and all Stieltjes orders}
\label{sec:hurwitz}

The new coefficient densities naturally supply the polynomial family
\begin{equation}\label{eq:Wkernel}
 W_{b,n}(v)=(1-v^2)^3(1+3v^2)^b(v+v^3)^{n-2},
 \qquad b\ge0,\ n\ge2.
\end{equation}
Multiplication by $v^k$ gives all coefficient moments in
\eqref{eq:coefficient-integral}. Every such polynomial and its first two
derivatives vanish at $v=\pm1$. This creates a finite endpoint calculus
starting with the third derivative, uniformly in $b,n,k$.

\subsection{The polynomial endpoint identity}
Let $P$ be any polynomial of degree $d$ and define
\[
 M_P(s;a)=\frac12\int_{-1}^1P(v)\zeta(s,a+v)\dd v,
 \qquad \mu_P=\avg{P(v)}.
\]
Use $\fall{s-1}{j+1}=(s-1)(s-2)\cdots(s-j-1)$.

\begin{proposition}[Finite Hurwitz endpoint formula]\label{prop:Hurwitz}
For $\Re a>1$, as an identity of meromorphic functions of $s$,
\begin{equation}\label{eq:Hurwitz-endpoint}
 M_P(s;a)=-\frac12\sum_{j=0}^d
 \frac{\eval{P^{(j)}(v)\zeta(s-j-1,a+v)}}{\fall{s-1}{j+1}}.
\end{equation}
Its only possible pole is at $s=1$, with residue $\mu_P$.
All other apparent singularities of the right side are removable after
summation. If $P,P',P''$ vanish at both endpoints, the sum starts at
$j=3$.
\end{proposition}
\begin{proof}
The classical argument derivative
$\partial_a\zeta(s-1,a)=-(s-1)\zeta(s,a)$
\cite{DLMF25} gives the first integration by parts. Repeat until
$P^{(d+1)}=0$. Each endpoint term has the minus sign shown in
\eqref{eq:Hurwitz-endpoint}. This proves the formula away from the
exceptional integers. On the compact interval and locally uniformly for
$\Re a>1$, the integrand is analytic in $s$ except for its simple
pole at one. Its residue integrates to $\mu_P$. Meromorphic continuation
therefore proves the endpoint identity and all stated cancellations.
\end{proof}

This proposition is the standard polynomial integration-by-parts
mechanism. The specialization to $W_{b,n}$ is what connects the newly
computed Gamma normal coefficients to spectral special functions.
No endpoint regularization in the variable $v$ is needed: the integrals
are ordinary compact-interval integrals. Laurent coefficients in $s$
are a different operation.

\subsection{Finite harmonic coefficients for every Stieltjes order}
For $j\ge0$ define $\eta_{j,\ell}$ by
\begin{equation}\label{eq:eta}
 \prod_{r=1}^j(1-\epsilon/r)^{-1}
 =\sum_{\ell\ge0}\eta_{j,\ell}\epsilon^\ell
 =\exp\left(\sum_{r\ge1}\frac{H_j^{(r)}}r\epsilon^r\right).
\end{equation}
The empty product at $j=0$ is one. In particular,
\[
 \eta_{j,0}=1,\qquad \eta_{j,1}=H_j,\qquad
 \eta_{j,2}=\frac{H_j^2+H_j^{(2)}}2.
\]
These are finite polynomials in generalized harmonic numbers at each
fixed $\ell$.

\begin{theorem}[All-order Stieltjes averages]\label{thm:Stieltjes}
For $q\ge0$ and $\Re a>1$,
\begin{equation}\label{eq:Stieltjes-endpoint}
\begin{split}
 \avg{P(v)\gamma_q(a+v)}
 ={}&\frac{(-1)^{q+1}q!}{2}
 \sum_{j=0}^d\frac{(-1)^j}{j!}\\[-2pt]
 &\quad\times\eval{P^{(j)}(v)
 \sum_{r=0}^{q+1}\eta_{j,q+1-r}
       \frac{\zeta^{(r)}(-j,a+v)}{r!}}.
\end{split}
\end{equation}
The formula is finite at every $q$ and applies in particular to every
polynomial $v^kW_{b,n}$.
\end{theorem}
\begin{proof}
At $s=1+\epsilon$,
\begin{equation}\label{eq:falling-laurent}
 \frac1{\fall{\epsilon}{j+1}}
 =\frac{(-1)^j}{j!\epsilon}
       \sum_{\ell\ge0}\eta_{j,\ell}\epsilon^\ell.
\end{equation}
Expand $\zeta(-j+\epsilon,a+v)$ in its order variable and insert
\eqref{eq:falling-laurent} into \eqref{eq:Hurwitz-endpoint}. Coefficient
$\epsilon^q$ is the right side of \eqref{eq:Stieltjes-endpoint}
multiplied by $(-1)^q/q!$. On the integral side the same coefficient is
$(-1)^q\avg{P\gamma_q}/q!$ by
\eqref{eq:stieltjes-convention}. Uniformity on compact parameter sets
justifies coefficient extraction under the integral. The coefficient of
$\epsilon^{-1}$ is the separately retained residue $\mu_P$; it must not
be discarded before multiplying the expansions.
\end{proof}

\begin{corollary}[Anchored primitives and argument derivatives]
\label{cor:Stieltjes-calculus}
For $a_0,a$ in the half-plane $\Re a>1$, with the integral taken along
any path in that half-plane,
\begin{equation}\label{eq:Stieltjes-primitive}
 \int_{a_0}^a\avg{P(v)\gamma_q(c+v)}\dd c
 =\left.\frac{(-1)^{q+1}}{q+1}\partial_s^{q+1}
       \bigl(M_P(s;a)-M_P(s;a_0)\bigr)\right|_{s=0}.
\end{equation}
For $r\ge1$, let $e_j^{[r]}$ be the elementary symmetric polynomial of
degree $j$ in $1,1/2,\ldots,1/r$. Then
\begin{equation}\label{eq:Stieltjes-argument-derivative}
\begin{split}
 \partial_a^r\avg{P(v)\gamma_q(a+v)}
 ={}&(-1)^{q+r}q!r!
 \sum_{j=0}^{\min(q,r)}e_j^{[r]}
 \frac{\left.\partial_s^{q-j}M_P(s;a)\right|_{s=r+1}}{(q-j)!}.
\end{split}
\end{equation}
Here the derivative $\partial_s$ on $M_P(r+1;a)$ is evaluated in its
first argument. Both formulas reduce to finite endpoint expressions by
Proposition~\ref{prop:Hurwitz} and its removable limits.
\end{corollary}
\begin{proof}
From $\partial_a\zeta(s,a)=-s\zeta(s+1,a)$ and the Laurent expansion,
\[
 \partial_a\zeta^{(q+1)}(0,a)
 =(-1)^{q+1}(q+1)\gamma_q(a).
\]
Integrate in $a$ and average against $P$ to obtain
\eqref{eq:Stieltjes-primitive}; the displayed subtraction fixes the
anchor exactly. For the derivative formula, use
\[
 \partial_a^r\zeta(s,a)=(-1)^r(s)_r\zeta(s+r,a),\qquad
 (1+\epsilon)_r=r!\sum_{j=0}^re_j^{[r]}\epsilon^j,
\]
and compare coefficients at $s=1+\epsilon$.
\end{proof}

For $q=0$, the primitive is the weighted log-Gamma difference because
$\zeta'(0,a)=\log\Gamma(a)-\tfrac12\log(2\pi)$
\cite{DLMF25}. Formula~\eqref{eq:Stieltjes-argument-derivative} at
$q=0$ recovers the weighted polygamma integrals, using
$\psi^{(r)}(a)=(-1)^{r+1}r!\zeta(r+1,a)$.

\subsection{Explicit digamma and Stieltjes evaluations}
\begin{corollary}[Four concrete evaluations]\label{cor:four-integrals}
All integrals below are ordinary integrals over $[-1,1]$:
\begin{align}
 \int_{-1}^1P_3(v)\psi(2+v)\dd v
 &=16\zeta'(-3)+12\zeta'(-5)+\frac27,
 \label{eq:I-psi-base}\\
 \int_{-1}^1P_3(v)(1+3v^2)\psi(2+v)\dd v
 &=64\zeta'(-3)+240\zeta'(-5)+48\zeta'(-7)
       +\frac{547}{1800},\label{eq:I-psi-mixed}\\
 \int_{-1}^1P_3(v)(v+v^3)\psi(2+v)\dd v
 &=\frac{32}{105}-\frac{84\zeta(5)}{\pi^4}
   +\frac{1575\zeta(7)}{2\pi^6}-\frac{2835\zeta(9)}{2\pi^8},
 \label{eq:I-psi-odd}\\
 \int_{-1}^1P_3(v)\gamma_1(2+v)\dd v
 &=8\zeta''(-3)+6\zeta''(-5)+\frac{88}{3}\zeta'(-3)
   +\frac{137}{5}\zeta'(-5)\notag\\
 &\hspace{1em}-\frac{16}{15}\log2+\frac{521}{900}.
 \label{eq:I-stieltjes}
\end{align}
\end{corollary}
\begin{proof}
At $a=2$ the endpoint arguments are $1$ and $3$. Use the exact shift
\begin{equation}\label{eq:zeta-three}
 \zeta(s,3)=\zeta(s)-1-2^{-s}.
\end{equation}
For $q=0$, Theorem~\ref{thm:Stieltjes} and $\gamma_0=-\psi$ give
\begin{equation}\label{eq:psi-endpoint}
 \int_{-1}^1P(v)\psi(2+v)\dd v
 =\sum_{j=0}^d\frac{(-1)^j}{j!}
 \eval{P^{(j)}(v)\bigl(\zeta'(-j,2+v)+H_j\zeta(-j,2+v)\bigr)}.
\end{equation}
The nonzero derivatives at $v=1$ needed for the three polynomials are
listed below; their derivatives at $-1$ follow from parity.
\begin{center}\small
\begin{tabular}{r|rrr}
\toprule
$j$ & $P_3^{(j)}(1)$ & $(P_3(1+3v^2))^{(j)}(1)$
    & $(P_3(v+v^3))^{(j)}(1)$\\
\midrule
3 & $-48$ & $-192$ & $-96$\\
4 & $-288$ & $-2304$ & $-1344$\\
5 & $-720$ & $-14400$ & $-10080$\\
6 & $-720$ & $-54720$ & $-50400$\\
7 & $0$ & $-120960$ & $-171360$\\
8 & $0$ & $-120960$ & $-362880$\\
9 & $0$ & $0$ & $-362880$\\
\bottomrule
\end{tabular}
\end{center}
Substituting these finite lists into \eqref{eq:psi-endpoint} and using
\eqref{eq:zeta-three} proves \eqref{eq:I-psi-base} and
\eqref{eq:I-psi-mixed}. The third polynomial first gives
\[
 -112\zeta'(-4)-140\zeta'(-6)-18\zeta'(-8)+\frac{32}{105}.
\]
The zeta functional equation implies
$\zeta'(-2r)=(-1)^r(2r)!\zeta(2r+1)/(2(2\pi)^{2r})$;
substitution gives \eqref{eq:I-psi-odd}.

Finally use $q=1$ in \eqref{eq:Stieltjes-endpoint}. Its inner sum is
\[
 \tfrac12\zeta''(-j,2+v)+H_j\zeta'(-j,2+v)
 +\tfrac12(H_j^2+H_j^{(2)})\zeta(-j,2+v).
\]
With the first derivative column and \eqref{eq:zeta-three}, collecting
terms yields \eqref{eq:I-stieltjes}. These calculations require no
numerical recognition of constants.
\end{proof}

The second and third integrals use precisely the first mixed square and
nonlinear cubic coefficient kernels, apart from their fixed scaling
constants. The theorem thus links the algebraic coefficients to explicit
zeta and Stieltjes periods. It does not assert that these particular
low-degree evaluations have no earlier appearance in the literature.

% END INCLUDED FILE: sections/07_hurwitz_stieltjes.tex

% BEGIN INCLUDED FILE: sections/08_polylogarithms.tex
\section{Polylogarithm sums and Lerch order derivatives}
\label{sec:polylog}

The beta moments in the mixed coefficient formula can be summed directly.
This gives ordinary polylogarithms, not only order derivatives or formal
spectral symbols. We then give a second family involving Lerch and
polylogarithm order derivatives.

\subsection{An exact polylogarithm sum for every square kernel}
Write
\begin{equation}\label{eq:Pbcoeff}
 P_b(v)=(1-v^2)^3(1+3v^2)^b
       =\sum_{j=0}^{b+3}d_{b,j}v^{2j},
 \qquad c_b=\frac1{18432b!24^b}.
\end{equation}
Then \eqref{eq:gb-density} gives
\begin{equation}\label{eq:gb-rational}
 g_b(2\ell)=-c_b\sum_{j=0}^{b+3}\frac{d_{b,j}}{2\ell+2j+1}.
\end{equation}
The elementary odd-part polylogarithm is
\[
 \chi_r(z)=\frac{\Li_r(z)-\Li_r(-z)}2
 =\sum_{\ell\ge0}\frac{z^{2\ell+1}}{(2\ell+1)^r},\qquad |z|<1.
\]
The defining polylogarithm series and its analytic conventions are
standard \cite{DLMF12}.

\begin{theorem}[All coefficient-weighted odd polylogarithm sums]
\label{thm:polylog-sum}
Let $b\ge0$, $m\ge1$, and $|z|<1$. Then
\begin{equation}\label{eq:polylog-sum}
 \sum_{\ell\ge0}\frac{g_b(2\ell)z^{2\ell+1}}{(2\ell+1)^m}
 =-c_b\sum_{j=0}^{b+3}d_{b,j}T_{m,j}(z),
\end{equation}
where $T_{m,0}(z)=\chi_{m+1}(z)$ and, for $j\ge1$,
\begin{equation}\label{eq:T-polylog}
\begin{split}
 T_{m,j}(z)={}&\sum_{r=1}^m
 \frac{(-1)^{m-r}}{(2j)^{m-r+1}}\chi_r(z)\\
 &+\frac{(-1)^m z^{-2j}}{(2j)^m}
 \left(\chi_1(z)-\sum_{r=0}^{j-1}\frac{z^{2r+1}}{2r+1}\right).
\end{split}
\end{equation}
The apparent negative powers at $z=0$ are removable.
\end{theorem}
\begin{proof}
Substitution of \eqref{eq:gb-rational} reduces the assertion to evaluating
\[
 T_{m,j}(z)=\sum_{\ell\ge0}
 \frac{z^{2\ell+1}}{(2\ell+1)^m(2\ell+2j+1)}.
\]
For $j=0$ this is $\chi_{m+1}$. For $j\ge1$, use the finite partial
fraction identity, with $x=2\ell+1$ and $c=2j$,
\begin{equation}\label{eq:partial-fraction}
 \frac1{x^m(x+c)}=
 \sum_{r=1}^m\frac{(-1)^{m-r}}{c^{m-r+1}x^r}
 +\frac{(-1)^m}{c^m(x+c)}.
\end{equation}
Sum the terms using absolute convergence. In the last term, shift the
odd index by $2j$, obtaining \eqref{eq:T-polylog}. The bracket starts
with $z^{2j+1}$, proving removability at zero.
\end{proof}

\begin{corollary}[Zeta and harmonic endpoint values]\label{cor:polylog-endpoint}
The limit at $z=1$ is given by the same finite sum
\eqref{eq:polylog-sum}, with
\begin{align}
 T_{m,0}(1)&=(1-2^{-m-1})\zeta(m+1),\label{eq:T0-one}\\
 T_{m,j}(1)&=\sum_{r=2}^m
 \frac{(-1)^{m-r}(1-2^{-r})\zeta(r)}{(2j)^{m-r+1}}
 +\frac{(-1)^{m+1}}{(2j)^m}
       \left(H_{2j}-\frac12H_j\right),\qquad j\ge1.
 \label{eq:Tj-one}
\end{align}
An empty sum in \eqref{eq:Tj-one} is zero.
\end{corollary}
\begin{proof}
Combine the two $\chi_1(z)$ terms in \eqref{eq:T-polylog} before
letting $z\to1$. Their coefficient is a constant multiple of
$1-z^{-2j}$; its product with the logarithmic singularity tends to zero.
For $r\ge2$, $\chi_r(1)=(1-2^{-r})\zeta(r)$, and the finite odd
harmonic sum is $H_{2j}-H_j/2$. The defining series for $T_{m,j}$
converges absolutely at one. So does the coefficient-weighted series;
indeed the beta formula gives $g_b(2\ell)=O(\ell^{-4})$ for fixed $b$.
These observations justify passage to the endpoint.
\end{proof}

For example, a base normalization consequence is
\begin{equation}\label{eq:example-odd-sum}
 \sum_{\ell\ge0}
 \frac1{(2\ell+1)^2(2\ell+3)(2\ell+5)(2\ell+7)}
 =\frac{\pi^2}{384}-\frac{17}{1080}.
\end{equation}
This example is an elementary partial-fraction evaluation. The uniform
formula \eqref{eq:polylog-sum} is useful because the coefficients
$g_b$ have already been identified with every mixed square sector,
including the previously unanswered $b>0$ family.

Termwise application of $z\partial_z$ lowers $m$, and integration
against $\dd z/z$ from zero raises $m$. Absolute convergence on compact
subdisks justifies both operations. Thus \eqref{eq:polylog-sum} also
provides anchored antiderivatives and derivative identities for the same
coefficient families, without introducing new undetermined constants.

\subsection{Lerch transforms and polylogarithm order derivatives}
For $|z|<1$ and $\Re a>1$, set
\[
 L_P(z,s;a)=\frac12\int_{-1}^1P(v)\Phi(z,s,a+v)\dd v,
 \qquad
 \Phi(z,s,a)=\sum_{n\ge0}\frac{z^n}{(n+a)^s}.
\]
The Lerch series is entire in $s$ on this disk, and
$\partial_a\Phi(z,s-1,a)=-(s-1)\Phi(z,s,a)$
\cite{DLMF14}. Hence the same proof as Proposition~\ref{prop:Hurwitz}
gives
\begin{equation}\label{eq:Lerch-endpoint}
 L_P(z,s;a)=-\frac12\sum_{j=0}^d
 \frac{\eval{P^{(j)}(v)\Phi(z,s-j-1,a+v)}}{\fall{s-1}{j+1}}.
\end{equation}
Unlike the Hurwitz case, \emph{all} the apparent poles of this expression,
including $s=1$, are removable after the finite sum is combined.

\begin{proposition}[Order-one Lerch average]\label{prop:Lerch-one}
For $|z|<1$ and $\Re a>1$,
\begin{equation}\label{eq:Lerch-one}
 L_P(z,1;a)=-\frac12\sum_{j=0}^d\frac{(-1)^j}{j!}
 \eval{P^{(j)}(v)
 \bigl(\partial_s\Phi(z,-j,a+v)+H_j\Phi(z,-j,a+v)\bigr)}.
\end{equation}
At $a=2$ this is a finite expression in $\Li_{-j}(z)$,
$\partial_s\Li_s(z)|_{s=-j}$, and elementary terms.
\end{proposition}
\begin{proof}
Expand \eqref{eq:Lerch-endpoint} at $s=1+\epsilon$ using
\eqref{eq:falling-laurent}. Entirety in $s$ forces the coefficient of
$\epsilon^{-1}$ to cancel. The constant coefficient is exactly
\eqref{eq:Lerch-one}. At the integer endpoints,
\begin{equation}\label{eq:Lerch-integer}
 \Phi(z,s,1)=z^{-1}\Li_s(z),\qquad
 \Phi(z,s,3)=z^{-3}\bigl(\Li_s(z)-z-z^22^{-s}\bigr).
\end{equation}
These follow by shifting the defining series and prove the final claim.
\end{proof}

For an entirely explicit implementation, put
$A_j(z)=P^{(j)}(1)z^{-3}-P^{(j)}(-1)z^{-1}$ and write
$\dot\Li_{-j}(z)=\partial_s\Li_s(z)|_{s=-j}$. The endpoint expression
inside the summation in \eqref{eq:Lerch-one}, at $a=2$, is
\begin{equation}\label{eq:Lerch-polylog-explicit}
\begin{split}
 &A_j(z)\bigl(\dot\Li_{-j}(z)+H_j\Li_{-j}(z)\bigr)\\
 &\qquad+P^{(j)}(1)
 \left(2^jz^{-1}\log2-H_j\bigl(z^{-2}+2^jz^{-1}\bigr)\right).
\end{split}
\end{equation}
The apparent singularities at $z=0$ cancel by
\eqref{eq:Lerch-integer}. Formula~\eqref{eq:Lerch-one} therefore
connects every polynomial $v^kW_{b,n}$ to a finite family of polylogarithm
order derivatives with no infinite reduction step.

The limits $z\to1$ and $s\to1$ should not be interchanged without
analysis: the Lerch transform is entire in $s$ for $|z|<1$, whereas its
Hurwitz counterpart can have the genuine residue $\mu_P$ at $s=1$.
The direct Hurwitz formula in Theorem~\ref{thm:Stieltjes} keeps that
residue and is the appropriate normalization for Stieltjes constants.

% END INCLUDED FILE: sections/08_polylogarithms.tex

% BEGIN INCLUDED FILE: sections/09_audit_research.tex
\section{Verification, source audit, and further research}
\label{sec:audit}

\subsection{What has actually been verified}
The exact verification program constructs the Bell polynomials by
\eqref{eq:Bell-recur} and reduces them by the independent terminating
algorithm of Proposition~\ref{prop:normal}. It then compares the resulting
rational coefficients with the beta formulas. This route does not use the
cubic variational recurrence to compute the normal coefficients.

The supplied recorded run has the following scope:
\begin{center}
\begin{tabular}{lr}
\toprule
Exact check & Count\\
\midrule
Full identities $A_m=N_m+\D H_m$, weights $0$--$30$ & 31\\
Functional support assertions in those normal forms & 2974\\
Cubic-sector coefficients, including prescribed parity zeros & 973\\
Quadratic-sector coefficients & 575\\
Affine-sector coefficients & 256\\
All-index symbolic density identities & 4\\
Independent beta versus double-sum comparisons & 216\\
Endpoint-vanishing assertions for the mixed kernels & 288\\
\bottomrule
\end{tabular}
\end{center}
The four symbolic identities involve an indeterminate $n$ where
appropriate; the other counts are finite regression checks.
A second exact script checks all four displayed spectral integrals, the four
low-$b$ rational coefficient families, twelve partial-fraction instances, and
the printed $H_8$ and $N_{11}$ against the certificates.
The exact run used Python 3.13.5 and SymPy 1.14.0. Rational normal-form
coefficients use Python's exact \texttt{Fraction} arithmetic. JSON
certificates include both $N_m$ and $H_m$ at weights 8, 11, 14, 20, 24,
and 30, and a separate file lists nonzero mixed examples.

The numerical program performs 30 independent diagnostics: eight
small-parameter comparisons of the compact integral with independently
reduced normal coefficients, three weighted digamma evaluations, nine
Hurwitz endpoint identities, nine polylogarithm sums, and one
first-Stieltjes integral. Most use 50 decimal digits; the Stieltjes
integral uses 35. The configured tolerances are $10^{-38}$ and
$10^{-28}$ respectively, and all tests pass. These are arbitrary-precision
floating-point diagnostics, not certified error intervals. A truncated
normal-series comparison is likewise a diagnostic, not a bound for its
uncomputed remainder.

No proof assistant was run. The article supplies analytic and algebraic
proofs, the exact script checks the finite rational identities, and the
numerical script tests the special-function normalizations. These are
three different forms of evidence and are not conflated.

\subsection{Corrections and safeguards from the source review}
Within the narrowly inspected mathematical core, no substantive error is
asserted. The normal decomposition, the finite-jet lemma, and the prior
quadratic recurrence are retained with attribution. The new density
solution is a continuation, not a purported correction to those results.

There is an editorial duplication in the inspected manuscript README:
the paragraph beginning with the complete 1,032-input first-letter $X$
family, its rank 1,008, and kernel dimension 24 appears twice
consecutively \cite{RepoMain}. The proposed edit is simply to remove the
second occurrence, after checking that the current version still has it.
No mathematical assertion in that paragraph is changed.

The following normalization safeguards are important for integration.
They are not claims that the inspected source made each error. First,
$\sigma_d\NF\ne\NF\sigma_d$ in general; the example after
Lemma~\ref{lem:finitejet} supplies an exact counterexample. Second,
$\sigma_3N_m$ is not the full normal polynomial once higher derivative
sectors occur, as \eqref{eq:N11} shows. Third, the polynomial primitive
$H_m$ is essential in ordinary definite integrals. Fourth, the Laurent
residue at $s=1$ must be retained when multiplying the endpoint expansion;
otherwise the generalized harmonic corrections in
\eqref{eq:Stieltjes-endpoint} can be lost. Finally, the finite rational
coefficient rank of a declared basis says nothing by itself about
arithmetic independence of the corresponding zeta periods.

The incoming archive names were visible, including packages concerning
Gamma normalizations, colored moments, and Stieltjes resonances. Their
contents were not exhaustively reconciled with this article. The
integration note therefore requests a focused overlap check rather than
claiming that this is the globally newest result in every one of those
packages. The repository remains unchanged by this delivery.

\subsection{An all-derivative-order conjecture}
\label{sec:conjecture}
The density formula has a striking extension which is supported by exact
checks but is not proved here. For $j\ge0$, put
\begin{equation}\label{eq:qj}
 q_j(v)=\frac{(1+v)^{j+1}-(v-1)^{j+1}}
               {2^{j+1}(j+1)!}.
\end{equation}
Then $q_0=1$, $q_1=v/2$, $q_2=h$, $q_3=k$, and
$q_j'=q_{j-1}/2$ for $j\ge1$.

\begin{conjecture}[Universal highest-derivative density]
\label{conj:universal}
Let $r\ge1$ and let
$M=u_0^a\prod_{j=1}^ru_j^{m_j}$ be functional with $m_r\ge2$ and
weight $W=a+\sum_{j=1}^r(j+1)m_j$. Then
\begin{equation}\label{eq:universal-conjecture}
\begin{split}
 [M]N_W={}&\frac{(-1)^rW!}
 {4^r r!(r+1)!\,a!\prod_{j=1}^rm_j!}\\
 &\times\avg{(1-v^2)^r q_r(v)^{m_r-2}
                   \prod_{j=1}^{r-1}q_j(v)^{m_j}}.
\end{split}
\end{equation}
\end{conjecture}

The cases $r=1,2,3$ follow from the affine formula,
Theorem~\ref{thm:quadratic}, and Theorem~\ref{thm:main}. The exact
exploratory program compares 438 nonzero coefficients with highest index
at least four through weight 24, and every one agrees. That finite
comparison does \emph{not} prove the conjecture, nor does it check all
absent monomials beyond the proved sectors.

In generating-function form the conjectured contribution of the largest
index $r$ is
\begin{equation}\label{eq:universal-increment}
 \frac{(-1)^r x_r^2}{4^r r!(r+1)!}
 \avg{(1-v^2)^r
   \exp\left(\sum_{j=1}^{r-1}x_jq_j(v)\right)\Erem(x_rq_r(v))},
 \qquad x_j=u_jt^{j+1},
\end{equation}
with the overall Appell factor $e^{u_0t}$. For $r=1$ this expression is
the affine correction $f_{\mathrm{aff}}(y)-1$ with $x_1=2y$.
The factor $(1-v^2)^r$, the derivative ladder for $q_j$, and the
normal-basis condition all point toward a general variational transport
proof. Such a proof is not supplied by the cubic calculation alone.

\subsection{Further research questions}
\paragraph{1. Prove or refute the universal density.}
Establish \eqref{eq:universal-conjecture} at arbitrary $r$, or exhibit a
first counterexample. A desirable proof would replace the order-three
polynomial identity by a uniform identity using $q_j'=q_{j-1}/2$ and the
$r$ endpoint zeros of $(1-v^2)^r$. The conjecture is specific enough for
independent symbolic or proof-assistant attack.

\paragraph{2. Find a closed endpoint formula, not only a terminating algorithm.}
The canonical $H_m$ is completely specified and verifiable at every
weight, but a compact generating representation analogous to
\eqref{eq:main} would make analytic integrations more transparent. It must
respect the constant normalization and distinguish a primitive of the
full normal form from one of a truncated projection.

\paragraph{3. Determine the arithmetic image of the new mixed kernels.}
For rational $a>1$, classify the finite Hurwitz-zeta derivative spaces
actually occupied by $\avg{v^kW_{b,n}(v)\gamma_q(a+v)}$. The endpoint
formula gives an explicit spanning family, but does not prove minimality
or numerical period independence. Distribution identities and parity may
yield substantially smaller exact families.

\paragraph{4. Reconcile Lerch confluence with the Hurwitz residue.}
Develop a uniform two-parameter formula near $(z,s)=(1,1)$ that separates
the emerging Hurwitz pole from the analytic Lerch remainder. This would
turn the present warning about order of limits into an exact confluence
theorem for every mixed kernel and its order derivatives.

\paragraph{5. Evaluate additional nonlinear polygamma combinations.}
The formulas here supply exact coefficient identities and retained
Gamma endpoints. A further task is to evaluate selected individual
integrals, for example those involving $\psi''(\psi''')^2$, in explicitly
declared convergent zeta-derivative or multiple-polylogarithm coordinates.
The high-derivative remainder in \eqref{eq:projected-endpoint} must remain
visible until it is independently reduced.

\paragraph{6. Study convergence beyond a fixed projection.}
Every fixed cubic projection is entire in its finite set of variables.
This does not establish convergence of a sum over all derivative indices
after substituting the infinite jet of digamma. A rigorous study should
separate formal coefficient identities from analytic convergence of the
full normal-form generating series.

\paragraph{7. Formalize the proof-carrying coefficient compiler.}
Formalize the terminating monomial reduction and the identity
$A_m=N_m+\D H_m$, then verify the rational certificate files. The
all-index density identities are a compact subsequent target because they
reduce to polynomial arithmetic with an indeterminate index.

\paragraph{8. Preserve the separate period problems.}
The present work does not resolve the general unequal-even-tail optimal
prefix problem, the Gaussian $S_6$ identity, or the revised $S_8$ candidate
in the inspected manuscript. Their exact branch conventions and lower
weight constants remain separate research tasks. No conclusion about
those problems follows from the coefficient formula proved here.

\subsection*{Conclusion}
A previously coupled infinite coefficient problem collapses to one entire
beta-density formula. The endpoint factor that proves the formula also
organizes exact Hurwitz, Stieltjes, Lerch, and ordinary polylogarithm
identities. The result is an all-weight theorem in a specified normal
basis, with retained primitives and reproducible checks, rather than a
finite table promoted to a conjectural theorem. The universal extension
has been deliberately left at its actual status: a precise and strongly
tested further question.

% END INCLUDED FILE: sections/09_audit_research.tex

% BEGIN INCLUDED FILE: references.tex
\clearpage
\begingroup
\small\raggedright
\begin{thebibliography}{99}
\addcontentsline{toc}{section}{References}

\bibitem{RepoMain}
ProveIt Contributors,
\emph{Polylogarithms and their Arithmetic Bridges}, manuscript README,
VladimirReshetnikov/ProveIt, inspected 11 October 2026.
\url{https://github.com/VladimirReshetnikov/ProveIt/tree/main/Analysis/Polylogarithms/docs/manuscript}.
File-level provenance is recorded in the accompanying manifest.

\bibitem{RepoQuadratic}
ProveIt research continuation,
\emph{A variational equation for all Gamma coefficients of derivative
order at most two}, section 04 of \emph{Harmonic Tail Gamma},
\path{Analysis/Polylogarithms/docs/reports/stieltjes-harmonic-resonance/continuations/harmonic-tail-gamma/sections/04_gamma_sectors.tex},
VladimirReshetnikov/ProveIt, inspected 11 October 2026.
Contains the normal-form construction, finite-jet Euler identity, and the
pure-$u_2$ recurrence used here.

\bibitem{RepoQuestions}
ProveIt research continuation,
\emph{Further research questions and concrete next targets}, section 11
of \emph{Affine Gamma Polarization},
\path{Analysis/Polylogarithms/docs/reports/stieltjes-harmonic-resonance/continuations/affine-gamma-polarization/sections/11_questions.tex},
VladimirReshetnikov/ProveIt, inspected 11 October 2026.
Question 1 is the coupled square-sector problem answered in
\eqref{eq:square-answer}.

\bibitem{GRR}
L. Guo, G. Regensburger, and M. Rosenkranz,
\emph{On integro-differential algebras},
Journal of Pure and Applied Algebra \textbf{218} (2014), 456--473.
\url{https://arxiv.org/abs/1212.0266}.
DOI: 10.1016/j.jpaa.2013.06.015.
The source continuation attributes the functional-monomial framework to
Proposition 4.10; the elementary decomposition needed here is reproved in
Proposition~\ref{prop:normal}.

\bibitem{DLMF5}
NIST Digital Library of Mathematical Functions,
Chapter 5, especially \S5.15, \emph{Polygamma Functions}.
\url{https://dlmf.nist.gov/5.15}.
Consulted for standard Gamma and polygamma conventions.

\bibitem{DLMF25}
NIST Digital Library of Mathematical Functions,
\S25.11, \emph{Hurwitz Zeta Function}.
\url{https://dlmf.nist.gov/25.11}.
In particular, the shift, argument-derivative, special-value, and
Gamma--zeta identities.

\bibitem{DLMF12}
NIST Digital Library of Mathematical Functions,
\S25.12, \emph{Polylogarithms}.
\url{https://dlmf.nist.gov/25.12}.

\bibitem{DLMF14}
NIST Digital Library of Mathematical Functions,
\S25.14, \emph{Lerch's Transcendent}.
\url{https://dlmf.nist.gov/25.14}.

\end{thebibliography}

\endgroup

% END INCLUDED FILE: references.tex

\end{document}
